% Source: book pp. 28--38 (PDF 42--52); solutions pp. 25--34.
% \axeqnum{N.M}: a displayed equation that Axler numbers with the shared
% item counter (2.20); use inside \[ ... \] (same definition in 2B, 2C).
\providecommand\axeqnum[1]{\refstepcounter{theorem}\tag*{\thetheorem}\label{ax:#1}}
\section{Span and Linear Independence}

We have been writing lists of numbers surrounded by parentheses, and we will
continue to do so for elements of $\F^n$; for example, $(2,-7,8)\in\F^3$.
However, now we need to consider lists of vectors (which may be elements of
$\F^n$ or of other vector spaces). To avoid confusion, we will usually write
lists of vectors without surrounding parentheses. For example,
$(4,1,6),(9,5,7)$ is a list of length two of vectors in $\R^3$.

\begin{notation}[List of vectors]\label{ax:2.1}
We will usually write lists of vectors without surrounding parentheses.
\end{notation}

\subsection{Linear Combinations and Span}

A sum of scalar multiples of the vectors in a list is called a linear
combination of the list. Here is the formal definition.

\begin{definition}[Linear combination]\label{ax:2.2}
A \emph{linear combination} of a list $v_1,\dots,v_m$ of vectors in $V$ is a
vector of the form
\[ a_1v_1+\cdots+a_mv_m, \]
where $a_1,\dots,a_m\in\F$.
\end{definition}

\begin{example}[Linear combinations in $\R^3$]\label{ax:2.3}\leavevmode
\begin{itemize}
\item $(17,-4,2)$ is a linear combination of $(2,1,-3),(1,-2,4)$, which is a
  list of length two of vectors in $\R^3$, because
  \[ (17,-4,2) = 6(2,1,-3)+5(1,-2,4). \]
\item $(17,-4,5)$ is not a linear combination of $(2,1,-3),(1,-2,4)$, which is
  a list of length two of vectors in $\R^3$, because there do not exist numbers
  $a_1,a_2\in\R$ such that
  \[ (17,-4,5) = a_1(2,1,-3)+a_2(1,-2,4). \]
  In other words, the system of equations
  \begin{align*}
    17 &= 2a_1+a_2\\
    -4 &= a_1-2a_2\\
     5 &= -3a_1+4a_2
  \end{align*}
  has no solutions (as you should verify).
\end{itemize}
\end{example}

\begin{definition}[Span]\label{ax:2.4}
The set of all linear combinations of a list of vectors $v_1,\dots,v_m$ in $V$
is called the \emph{span} of $v_1,\dots,v_m$, denoted by
$\spn(v_1,\dots,v_m)$. In other words,
\[ \spn(v_1,\dots,v_m) = \{a_1v_1+\cdots+a_mv_m : a_1,\dots,a_m\in\F\}. \]
The span of the empty list $(\,)$ is defined to be $\{0\}$.
\end{definition}

\begin{example}[Span]\label{ax:2.5}
The previous example shows that in $\F^3$,
\begin{itemize}
\item $(17,-4,2)\in\spn\bigl((2,1,-3),(1,-2,4)\bigr)$;
\item $(17,-4,5)\notin\spn\bigl((2,1,-3),(1,-2,4)\bigr)$.
\end{itemize}
\end{example}

\begin{result}[Span is the smallest containing subspace]\label{ax:2.6}
The span of a list of vectors in $V$ is the smallest subspace of $V$
containing all vectors in the list.
\end{result}

\begin{proof}
Suppose $v_1,\dots,v_m$ is a list of vectors in $V$.

\marginnote[2em]{Some mathematicians use the terminology \textbf{linear span},
which means the same as span.}%
First we show that $\spn(v_1,\dots,v_m)$ is a subspace of $V$. The additive
identity is in $\spn(v_1,\dots,v_m)$ because
\[ 0 = 0v_1+\cdots+0v_m. \]
Also, $\spn(v_1,\dots,v_m)$ is closed under addition because
\[ (a_1v_1+\cdots+a_mv_m)+(c_1v_1+\cdots+c_mv_m)
   = (a_1+c_1)v_1+\cdots+(a_m+c_m)v_m. \]
Furthermore, $\spn(v_1,\dots,v_m)$ is closed under scalar multiplication
because
\[ \lambda(a_1v_1+\cdots+a_mv_m) = \lambda a_1v_1+\cdots+\lambda a_mv_m. \]
Thus $\spn(v_1,\dots,v_m)$ is a subspace of $V$ (by \ref{ax:1.34}).

Each $v_k$ is a linear combination of $v_1,\dots,v_m$ (to show this, set
$a_k=1$ and let the other $a$'s in \ref{ax:2.2} equal $0$). Thus
$\spn(v_1,\dots,v_m)$ contains each $v_k$. Conversely, because subspaces are
closed under scalar multiplication and addition, every subspace of $V$ that
contains each $v_k$ contains $\spn(v_1,\dots,v_m)$. Thus $\spn(v_1,\dots,v_m)$
is the smallest subspace of $V$ containing all the vectors $v_1,\dots,v_m$.
\end{proof}

\begin{definition}[Spans]\label{ax:2.7}
If $\spn(v_1,\dots,v_m)$ equals $V$, we say that the list $v_1,\dots,v_m$
\emph{spans} $V$.
\end{definition}

\begin{example}[A list that spans $\F^n$]\label{ax:2.8}
Suppose $n$ is a positive integer. We want to show that
\[ (1,0,\dots,0),(0,1,0,\dots,0),\dots,(0,\dots,0,1) \]
spans $\F^n$. Here the $k^{\text{th}}$ vector in the list above has $1$ in the
$k^{\text{th}}$ slot and $0$ in all other slots.

Suppose $(x_1,\dots,x_n)\in\F^n$. Then
\[ (x_1,\dots,x_n) = x_1(1,0,\dots,0)+x_2(0,1,0,\dots,0)+\cdots
   +x_n(0,\dots,0,1). \]
Thus $(x_1,\dots,x_n)\in
\spn\bigl((1,0,\dots,0),(0,1,0,\dots,0),\dots,(0,\dots,0,1)\bigr)$, as
desired.
\end{example}

Now we can make one of the key definitions in linear algebra.

\begin{definition}[Finite-dimensional vector space]\label{ax:2.9}
A vector space is called \emph{finite-dimensional} if some list of vectors in
it spans the space.
\end{definition}

\marginnote{Recall that by definition every list has finite length.}%
Example~\ref{ax:2.8} above shows that $\F^n$ is a finite-dimensional vector
space for every positive integer $n$.

The definition of a polynomial is no doubt already familiar to you.

\begin{definition}[Polynomial, $\Poly(\F)$]\label{ax:2.10}\leavevmode
\begin{itemize}
\item A function $p\colon\F\to\F$ is called a \emph{polynomial} with
  coefficients in $\F$ if there exist $a_0,\dots,a_m\in\F$ such that
  \[ p(z) = a_0+a_1z+a_2z^2+\cdots+a_mz^m \]
  for all $z\in\F$.
\item $\Poly(\F)$ is the set of all polynomials with coefficients in $\F$.
\end{itemize}
\end{definition}

With the usual operations of addition and scalar multiplication, $\Poly(\F)$ is
a vector space over $\F$, as you should verify. Hence $\Poly(\F)$ is a subspace
of $\F^{\F}$, the vector space of functions from $\F$ to $\F$.

If a polynomial (thought of as a function from $\F$ to $\F$) is represented by
two sets of coefficients, then subtracting one representation of the
polynomial from the other produces a polynomial that is identically zero as a
function on $\F$ and hence has all zero coefficients (if you are unfamiliar
with this fact, just believe it for now; we will prove it later---see
\ref{ax:4.8}). \textbf{Conclusion:} the coefficients of a polynomial are
uniquely determined by the polynomial. Thus the next definition uniquely
defines the degree of a polynomial.

\begin{definition}[Degree of a polynomial, $\deg p$]\label{ax:2.11}\leavevmode
\begin{itemize}
\item A polynomial $p\in\Poly(\F)$ is said to have \emph{degree} $m$ if there
  exist scalars $a_0,a_1,\dots,a_m\in\F$ with $a_m\ne0$ such that for every
  $z\in\F$, we have
  \[ p(z) = a_0+a_1z+\cdots+a_mz^m. \]
\item The polynomial that is identically $0$ is said to have degree $-\infty$.
\item The degree of a polynomial $p$ is denoted by $\deg p$.
\end{itemize}
\end{definition}

In the next definition, we use the convention that $-\infty<m$, which means
that the polynomial $0$ is in $\Poly_m(\F)$.

\begin{notation}[$\Poly_m(\F)$]\label{ax:2.12}
For $m$ a nonnegative integer, $\Poly_m(\F)$ denotes the set of all polynomials
with coefficients in $\F$ and degree at most $m$.
\end{notation}

If $m$ is a nonnegative integer, then $\Poly_m(\F)=\spn(1,z,\dots,z^m)$ [here
we slightly abuse notation by letting $z^k$ denote a function]. Thus
$\Poly_m(\F)$ is a finite-dimensional vector space for each nonnegative integer
$m$.

\begin{definition}[Infinite-dimensional vector space]\label{ax:2.13}
A vector space is called \emph{infinite-dimensional} if it is not
finite-dimensional.
\end{definition}

\begin{example}[$\Poly(\F)$ is infinite-dimensional]\label{ax:2.14}
Consider any list of elements of $\Poly(\F)$. Let $m$ denote the highest degree
of the polynomials in this list. Then every polynomial in the span of this list
has degree at most $m$. Thus $z^{m+1}$ is not in the span of our list. Hence no
list spans $\Poly(\F)$. Thus $\Poly(\F)$ is infinite-dimensional.
\end{example}

\subsection{Linear Independence}

Suppose $v_1,\dots,v_m\in V$ and $v\in\spn(v_1,\dots,v_m)$. By the definition
of span, there exist $a_1,\dots,a_m\in\F$ such that
\[ v = a_1v_1+\cdots+a_mv_m. \]
Consider the question of whether the choice of scalars in the equation above is
unique. Suppose $c_1,\dots,c_m$ is another set of scalars such that
\[ v = c_1v_1+\cdots+c_mv_m. \]
Subtracting the last two equations, we have
\[ 0 = (a_1-c_1)v_1+\cdots+(a_m-c_m)v_m. \]
Thus we have written $0$ as a linear combination of $(v_1,\dots,v_m)$. If the
only way to do this is by using $0$ for all the scalars in the linear
combination, then each $a_k-c_k$ equals $0$, which means that each $a_k$ equals
$c_k$ (and thus the choice of scalars was indeed unique). This situation is so
important that we give it a special name---linear independence---which we now
define.

\begin{definition}[Linearly independent]\label{ax:2.15}\leavevmode
\begin{itemize}
\item A list $v_1,\dots,v_m$ of vectors in $V$ is called \emph{linearly
  independent} if the only choice of $a_1,\dots,a_m\in\F$ that makes
  \[ a_1v_1+\cdots+a_mv_m = 0 \]
  is $a_1=\cdots=a_m=0$.
\item The empty list $(\,)$ is also declared to be linearly independent.
\end{itemize}
\end{definition}

The reasoning above shows that $v_1,\dots,v_m$ is linearly independent if and
only if each vector in $\spn(v_1,\dots,v_m)$ has only one representation as a
linear combination of $v_1,\dots,v_m$.

\begin{example}[Linearly independent lists]\label{ax:2.16}\leavevmode
\begin{parts}
\item To see that the list $(1,0,0,0),(0,1,0,0),(0,0,1,0)$ is linearly
  independent in $\F^4$, suppose $a_1,a_2,a_3\in\F$ and
  \[ a_1(1,0,0,0)+a_2(0,1,0,0)+a_3(0,0,1,0) = (0,0,0,0). \]
  Thus
  \[ (a_1,a_2,a_3,0) = (0,0,0,0). \]
  Hence $a_1=a_2=a_3=0$. Thus the list $(1,0,0,0),(0,1,0,0),(0,0,1,0)$ is
  linearly independent in $\F^4$.
\item Suppose $m$ is a nonnegative integer. To see that the list
  $1,z,\dots,z^m$ is linearly independent in $\Poly(\F)$, suppose
  $a_0,a_1,\dots,a_m\in\F$ and
  \[ a_0+a_1z+\cdots+a_mz^m = 0, \]
  where we think of both sides as elements of $\Poly(\F)$. Then
  \[ a_0+a_1z+\cdots+a_mz^m = 0 \]
  for all $z\in\F$. As discussed earlier (and as follows from \ref{ax:4.8}),
  this implies that $a_0=a_1=\cdots=a_m=0$. Thus $1,z,\dots,z^m$ is a linearly
  independent list in $\Poly(\F)$.
\item A list of length one in a vector space is linearly independent if and
  only if the vector in the list is not $0$.
\item A list of length two in a vector space is linearly independent if and
  only if neither of the two vectors in the list is a scalar multiple of the
  other.
\end{parts}
\end{example}

If some vectors are removed from a linearly independent list, the remaining
list is also linearly independent, as you should verify.

\begin{definition}[Linearly dependent]\label{ax:2.17}\leavevmode
\begin{itemize}
\item A list of vectors in $V$ is called \emph{linearly dependent} if it is not
  linearly independent.
\item In other words, a list $v_1,\dots,v_m$ of vectors in $V$ is linearly
  dependent if there exist $a_1,\dots,a_m\in\F$, not all $0$, such that
  $a_1v_1+\cdots+a_mv_m=0$.
\end{itemize}
\end{definition}

\begin{example}[Linearly dependent lists]\label{ax:2.18}\leavevmode
\begin{itemize}
\item $(2,3,1),(1,-1,2),(7,3,8)$ is linearly dependent in $\F^3$ because
  \[ 2(2,3,1)+3(1,-1,2)+(-1)(7,3,8) = (0,0,0). \]
\item The list $(2,3,1),(1,-1,2),(7,3,c)$ is linearly dependent in $\F^3$ if
  and only if $c=8$, as you should verify.
\item If some vector in a list of vectors in $V$ is a linear combination of the
  other vectors, then the list is linearly dependent. (Proof: After writing one
  vector in the list as equal to a linear combination of the other vectors,
  move that vector to the other side of the equation, where it will be
  multiplied by $-1$.)
\item Every list of vectors in $V$ containing the $0$ vector is linearly
  dependent. (This is a special case of the previous bullet point.)
\end{itemize}
\end{example}

The next lemma is a terrific tool. It states that given a linearly dependent
list of vectors, one of the vectors is in the span of the previous ones.
Furthermore, we can throw out that vector without changing the span of the
original list.

\begin{result}[Linear dependence lemma]\label{ax:2.19}
Suppose $v_1,\dots,v_m$ is a linearly dependent list in $V$. Then there exists
$k\in\{1,2,\dots,m\}$ such that
\[ v_k\in\spn(v_1,\dots,v_{k-1}). \]
Furthermore, if $k$ satisfies the condition above and the $k^{\text{th}}$ term
is removed from $v_1,\dots,v_m$, then the span of the remaining list equals
$\spn(v_1,\dots,v_m)$.
\end{result}

\begin{proof}
Because the list $v_1,\dots,v_m$ is linearly dependent, there exist numbers
$a_1,\dots,a_m\in\F$, not all $0$, such that
\[ a_1v_1+\cdots+a_mv_m = 0. \]
Let $k$ be the largest element of $\{1,\dots,m\}$ such that $a_k\ne0$. Then
\[ v_k = -\frac{a_1}{a_k}v_1-\cdots-\frac{a_{k-1}}{a_k}v_{k-1}, \]
which proves that $v_k\in\spn(v_1,\dots,v_{k-1})$, as desired.

Now suppose $k$ is any element of $\{1,\dots,m\}$ such that
$v_k\in\spn(v_1,\dots,v_{k-1})$. Let $b_1,\dots,b_{k-1}\in\F$ be such that
\[ v_k = b_1v_1+\cdots+b_{k-1}v_{k-1}. \axeqnum{2.20} \]
Suppose $u\in\spn(v_1,\dots,v_m)$. Then there exist $c_1,\dots,c_m\in\F$ such
that
\[ u = c_1v_1+\cdots+c_mv_m. \]
In the equation above, we can replace $v_k$ with the right side of
\ref{ax:2.20}, which shows that $u$ is in the span of the list obtained by
removing the $k^{\text{th}}$ term from $v_1,\dots,v_m$. Thus removing the
$k^{\text{th}}$ term of the list $v_1,\dots,v_m$ does not change the span of the
list.
\end{proof}

If $k=1$ in the linear dependence lemma, then $v_k\in\spn(v_1,\dots,v_{k-1})$
means that $v_1=0$, because $\spn(\,)=\{0\}$. Note also that parts of the proof
of the linear dependence lemma need to be modified if $k=1$. In general, the
proofs in the rest of the book will not call attention to special cases that
must be considered involving lists of length $0$, the subspace $\{0\}$, or
other trivial cases for which the result is true but needs a slightly
different proof. Be sure to check these special cases yourself.

\begin{example}[Smallest $k$ in linear dependence lemma]\label{ax:2.21}
Consider the list
\[ (1,2,3),(6,5,4),(15,16,17),(8,9,7) \]
in $\R^3$. This list of length four is linearly dependent, as we will soon see.
Thus the linear dependence lemma implies that there exists $k\in\{1,2,3,4\}$
such that the $k^{\text{th}}$ vector in this list is a linear combination of
the previous vectors in the list. Let's see how to find the smallest value of
$k$ that works.

Taking $k=1$ in the linear dependence lemma works if and only if the first
vector in the list equals $0$. Because $(1,2,3)$ is not the $0$ vector, we
cannot take $k=1$ for this list.

Taking $k=2$ in the linear dependence lemma works if and only if the second
vector in the list is a scalar multiple of the first vector. However, there
does not exist $c\in\R$ such that $(6,5,4)=c(1,2,3)$. Thus we cannot take
$k=2$ for this list.

Taking $k=3$ in the linear dependence lemma works if and only if the third
vector in the list is a linear combination of the first two vectors. Thus for
the list in this example, we want to know whether there exist $a,b\in\R$ such
that
\[ (15,16,17) = a(1,2,3)+b(6,5,4). \]
The equation above is equivalent to a system of three linear equations in the
two unknowns $a,b$. Using Gaussian elimination or appropriate software, we find
that $a=3$, $b=2$ is a solution of the equation above, as you can verify. Thus
for the list in this example, taking $k=3$ is the smallest value of $k$ that
works in the linear dependence lemma.
\end{example}

Now we come to a key result. It says that no linearly independent list in $V$
is longer than a spanning list in $V$.

\begin{result}[Length of linearly independent list $\le$ length of spanning
list]\label{ax:2.22}
In a finite-dimensional vector space, the length of every linearly independent
list of vectors is less than or equal to the length of every spanning list of
vectors.
\end{result}

\begin{proof}
Suppose that $u_1,\dots,u_m$ is linearly independent in $V$. Suppose also that
$w_1,\dots,w_n$ spans $V$. We need to prove that $m\le n$. We do so through the
process described below with $m$ steps; note that in each step we add one of
the $u$'s and remove one of the $w$'s.
\begin{description}[style=nextline,leftmargin=1em,itemsep=2pt,parsep=3pt]
\item[Step 1] Let $B$ be the list $w_1,\dots,w_n$, which spans $V$. Adjoining
  $u_1$ at the beginning of this list produces a linearly dependent list
  (because $u_1$ can be written as a linear combination of $w_1,\dots,w_n$). In
  other words, the list
  \[ u_1,w_1,\dots,w_n \]
  is linearly dependent.

  Thus by the linear dependence lemma (\ref{ax:2.19}), one of the vectors in
  the list above is a linear combination of the previous vectors in the list.
  We know that $u_1\ne0$ because the list $u_1,\dots,u_m$ is linearly
  independent. Thus $u_1$ is not in the span of the previous vectors in the
  list above (because $u_1$ is not in $\{0\}$, which is the span of the empty
  list). Hence the linear dependence lemma implies that we can remove one of
  the $w$'s so that the new list $B$ (of length $n$) consisting of $u_1$ and
  the remaining $w$'s spans $V$.
\item[Step $\bm{k}$, for $\bm{k=2},\dots,\bm{m}$] The list $B$ (of length $n$) from
  step $k-1$ spans $V$. In particular, $u_k$ is in the span of the list $B$.
  Thus the list of length $(n+1)$ obtained by adjoining $u_k$ to $B$, placing
  it just after $u_1,\dots,u_{k-1}$, is linearly dependent. By the linear
  dependence lemma (\ref{ax:2.19}), one of the vectors in this list is in the
  span of the previous ones, and because $u_1,\dots,u_k$ is linearly
  independent, this vector cannot be one of the $u$'s.

  Hence there still must be at least one remaining $w$ at this step. We can
  remove from our new list (after adjoining $u_k$ in the proper place) a $w$
  that is a linear combination of the previous vectors in the list, so that the
  new list $B$ (of length $n$) consisting of $u_1,\dots,u_k$ and the remaining
  $w$'s spans $V$.
\end{description}

After step $m$, we have added all the $u$'s and the process stops. At each step
as we add a $u$ to $B$, the linear dependence lemma implies that there is some
$w$ to remove. Thus there are at least as many $w$'s as $u$'s.
\end{proof}

The next two examples show how the result above can be used to show, without
any computations, that certain lists are not linearly independent and that
certain lists do not span a given vector space.

\begin{example}[No list of length $4$ is linearly independent in
$\R^3$]\label{ax:2.23}
The list $(1,0,0),(0,1,0),(0,0,1)$, which has length three, spans $\R^3$. Thus
no list of length larger than three is linearly independent in $\R^3$.

For example, we now know that $(1,2,3),(4,5,8),(9,6,7),(-3,2,8)$, which is a
list of length four, is not linearly independent in $\R^3$.
\end{example}

\begin{example}[No list of length $3$ spans $\R^4$]\label{ax:2.24}
The list $(1,0,0,0),(0,1,0,0),(0,0,1,0),(0,0,0,1)$, which has length four, is
linearly independent in $\R^4$. Thus no list of length less than four spans
$\R^4$.

For example, we now know that $(1,2,3,-5),(4,5,8,3),(9,6,7,-1)$, which is a
list of length three, does not span $\R^4$.
\end{example}

Our intuition suggests that every subspace of a finite-dimensional vector space
should also be finite-dimensional. We now prove that this intuition is
correct.

\begin{result}[Finite-dimensional subspaces]\label{ax:2.25}
Every subspace of a finite-dimensional vector space is finite-dimensional.
\end{result}

\begin{proof}
Suppose $V$ is finite-dimensional and $U$ is a subspace of $V$. We need to
prove that $U$ is finite-dimensional. We do this through the following
multistep construction.
\begin{description}[style=nextline,leftmargin=1em,itemsep=2pt,parsep=3pt]
\item[Step 1] If $U=\{0\}$, then $U$ is finite-dimensional and we are done. If
  $U\ne\{0\}$, then choose a nonzero vector $u_1\in U$.
\item[Step $\bm{k}$] If $U=\spn(u_1,\dots,u_{k-1})$, then $U$ is
  finite-dimensional and we are done. If $U\ne\spn(u_1,\dots,u_{k-1})$, then
  choose a vector $u_k\in U$ such that
  \[ u_k\notin\spn(u_1,\dots,u_{k-1}). \]
\end{description}

After each step, as long as the process continues, we have constructed a list
of vectors such that no vector in this list is in the span of the previous
vectors. Thus after each step we have constructed a linearly independent list,
by the linear dependence lemma (\ref{ax:2.19}). This linearly independent list
cannot be longer than any spanning list of $V$ (by \ref{ax:2.22}). Thus the
process eventually terminates, which means that $U$ is finite-dimensional.
\end{proof}

% ------------------------------------------------------------------ exercises
\begin{exc}

\begin{exercise}
Find a list of four distinct vectors in $\F^3$ whose span equals
\[ \{(x,y,z)\in\F^3 : x+y+z=0\}. \]
\end{exercise}
\begin{solution}
Clearly, the span of the list involves two vectors, since we may rewrite $z$ in
terms of $x$ and $y$ as $z=-x-y$. Thus, we may choose the following two
vectors:
\[ v_1=(1,0,-1), \qquad v_2=(0,1,-1). \]
To find two more vectors, we may pick any two distinct linear combinations of
$v_1$ and $v_2$. For example, we may select
\[ v_3=v_1+v_2=(1,1,-2), \qquad v_4=2v_1+3v_2=(2,3,-5). \]
Then
\[ \spn(v_1,v_2,v_3,v_4)=\spn(v_1,v_2)=\{(x,y,z)\in\F^3 \mid x+y+z=0\}.
   \qedhere \]
\end{solution}

\begin{exercise}
Prove or give a counterexample: If $v_1,v_2,v_3,v_4$ spans $V$, then the list
\[ v_1-v_2,\,v_2-v_3,\,v_3-v_4,\,v_4 \]
also spans $V$.
\end{exercise}
\begin{solution}
Let $U=\{v_1,v_2,v_3,v_4\}$ and $W=\{v_1-v_2,v_2-v_3,v_3-v_4,v_4\}$. Clearly,
$v_4\in\spn U$. Then we may write
\begin{align*}
v_3 &= (v_3-v_4)+v_4\in\spn W,\\
v_2 &= (v_2-v_3)+v_3\in\spn W,\\
v_1 &= (v_1-v_2)+v_2\in\spn W.
\end{align*}
Thus, $v_1,v_2,v_3,v_4\in\spn W$, and so $V=\spn U\subseteq\spn W$. Since $W$
is a list of vectors in $V$, we have $\spn W\subseteq V$. Therefore,
$\spn W=V$, and so $W$ spans $V$.
\end{solution}

\begin{exercise}
Suppose $v_1,\dots,v_m$ is a list of vectors in $V$. For $k\in\{1,\dots,m\}$,
let
\[ w_k = v_1+\cdots+v_k. \]
Show that $\spn(v_1,\dots,v_m)=\spn(w_1,\dots,w_m)$.
\end{exercise}
\begin{solution}
This is a generalization of Exercise 2A2 and a direct consequence of the fact
that $v_k=w_k-w_{k-1}$ for $k\in\{2,\dots,m\}$ and $v_1=w_1$. Since each $v_k$
is a linear combination of $w_1,\dots,w_k$, we have
$\spn(v_1,\dots,v_m)\subseteq\spn(w_1,\dots,w_m)$. On the other hand, each
$w_k$ is a linear combination of $v_1,\dots,v_k$, so
$\spn(w_1,\dots,w_m)\subseteq\spn(v_1,\dots,v_m)$. Therefore,
$\spn(v_1,\dots,v_m)=\spn(w_1,\dots,w_m)$.
\end{solution}

\begin{exercise}\leavevmode
\begin{parts}
\item Show that a list of length one in a vector space is linearly independent
  if and only if the vector in the list is not $0$.
\item Show that a list of length two in a vector space is linearly independent
  if and only if neither of the two vectors in the list is a scalar multiple of
  the other.
\end{parts}
\end{exercise}
\begin{solution}\leavevmode
\begin{parts}
\item Let $v$ be the vector in the list. If $v=0$, then the list is linearly
  dependent, since $0\cdot v=0$. If $v\ne0$, then the only way to write $0$ as
  a linear combination of $v$ is to take the scalar multiple to be $0$. Thus,
  the list is linearly independent.
\item Let $v_1$ and $v_2$ be the vectors in the list. If $v_1$ is a scalar
  multiple of $v_2$, then there exists a scalar $\lambda$ such that
  $v_1=\lambda v_2$. Then we may write
  \[ 0 = v_1-\lambda v_2, \]
  which shows that the list is linearly dependent. Conversely, if the list is
  linearly dependent, then there exist scalars $\alpha$ and $\beta$, not both
  zero, such that
  \[ \alpha v_1+\beta v_2 = 0. \]
  If $\alpha\ne0$, then
  \[ v_1 = -\frac{\beta}{\alpha}v_2. \]
  Similarly, if $\beta\ne0$, then
  \[ v_2 = -\frac{\alpha}{\beta}v_1. \]
  Thus, the list is linearly independent if and only if neither vector is a
  scalar multiple of the other.
\end{parts}
\end{solution}

\begin{exercise}
Find a number $t$ such that
\[ (3,1,4),(2,-3,5),(5,9,t) \]
is not linearly independent in $\R^3$.
\end{exercise}
\begin{solution}
To say that the list is not linearly independent is to say that there exist
scalars $\alpha,\beta,\gamma$, not all zero, such that
\[ \alpha(3,1,4)+\beta(2,-3,5)+\gamma(5,9,t)=(0,0,0). \]
This is equivalent to the following system of equations:
\begin{align*}
0 &= 3\alpha+2\beta+5\gamma,\\
0 &= \alpha-3\beta+9\gamma,\\
0 &= 4\alpha+5\beta+t\gamma.
\end{align*}
From the second equation, we obtain $\alpha=3\beta-9\gamma$. Substituting this
into the first equation, we have
\begin{align*}
0 &= 3(3\beta-9\gamma)+2\beta+5\gamma\\
  &= 11\beta-22\gamma.
\end{align*}
Thus, $\beta=2\gamma$. Substituting this into the expression for $\alpha$, we
have $\alpha=3(2\gamma)-9\gamma=-3\gamma$. Substituting these values into the
third equation, we have
\begin{align*}
0 &= 4(-3\gamma)+5(2\gamma)+t\gamma\\
  &= -12\gamma+10\gamma+t\gamma\\
  &= (t-2)\gamma.
\end{align*}
Since $\gamma\ne0$, we must have $t-2=0$, or $t=2$. Then the list is not
linearly independent when $t=2$.
\end{solution}

\begin{exercise}
Show that the list $(2,3,1),(1,-1,2),(7,3,c)$ is linearly dependent in $\F^3$
if and only if $c=8$.
\end{exercise}
\begin{solution}
We may proceed similarly to Exercise 2A5 to obtain the following system of
equations:
\begin{align*}
0 &= 2\alpha+\beta+7\gamma,\\
0 &= 3\alpha-\beta+3\gamma,\\
0 &= \alpha+2\beta+c\gamma.
\end{align*}
The second equation gives us $\beta=3\alpha+3\gamma$. Substituting this into
the first equation, we have
\begin{align*}
0 &= 2\alpha+(3\alpha+3\gamma)+7\gamma\\
  &= 5\alpha+10\gamma.
\end{align*}
Thus, $\alpha=-2\gamma$. Substituting this into the expression for $\beta$, we
have $\beta=3(-2\gamma)+3\gamma=-3\gamma$. Substituting these values into the
third equation, we have
\begin{align*}
0 &= (-2\gamma)+2(-3\gamma)+c\gamma\\
  &= -2\gamma-6\gamma+c\gamma\\
  &= (c-8)\gamma.
\end{align*}
Since $\gamma\ne0$, we must have $c-8=0$, or $c=8$. Then the list is linearly
dependent when $c=8$.

Conversely, if $c\ne8$, then $(c-8)\gamma=0$ implies $\gamma=0$. Then
$\alpha=-2\gamma=0$ and $\beta=-3\gamma=0$, so the only solution to the system
of equations is $\alpha=\beta=\gamma=0$. Thus, the list is linearly
independent when $c\ne8$.
\end{solution}

\begin{exercise}\leavevmode
\begin{parts}
\item Show that if we think of $\C$ as a vector space over $\R$, then the list
  $1+i,1-i$ is linearly independent.
\item Show that if we think of $\C$ as a vector space over $\C$, then the list
  $1+i,1-i$ is linearly dependent.
\end{parts}
\end{exercise}
\begin{solution}\leavevmode
\begin{parts}
\item Suppose $\alpha,\beta\in\R$ such that
  \[ \alpha(1+i)+\beta(1-i)=0. \]
  Then we have
  \[ (\alpha+\beta)+(\alpha-\beta)i=0. \]
  Since $0=0+0i$, we must have $\alpha+\beta=0$ and $\alpha-\beta=0$. Solving
  this system of equations, we find that $\alpha=\beta=0$. Thus, the list is
  linearly independent.
\item Suppose $a+bi,c+di\in\C$ such that
  \[ (a+bi)(1+i)+(c+di)(1-i)=0. \]
  Then we have
  \[ (a-b+c+d)+(a+b-c+d)i=0. \]
  Then we have the system of equations
  \begin{align*}
  0 &= a-b+c+d,\\
  0 &= a+b-c+d.
  \end{align*}
  Adding the two equations results in $0=2a+2d$, or $a=-d$. Substituting this
  into the first equation, we obtain $b=c$. Then we have infinitely many
  nontrivial solutions to the system of equations, such as $a=1$, $b=0$,
  $c=0$, $d=-1$. Thus, the list is linearly dependent.
\end{parts}
\end{solution}

\begin{exercise}
Suppose $v_1,v_2,v_3,v_4$ is linearly independent in $V$. Prove that the list
\[ v_1-v_2,\,v_2-v_3,\,v_3-v_4,\,v_4 \]
is also linearly independent.
\end{exercise}
\begin{solution}
Let $\alpha,\beta,\gamma,\delta\in\F$ such that
\[ \alpha(v_1-v_2)+\beta(v_2-v_3)+\gamma(v_3-v_4)+\delta v_4=0. \]
Then we have
\[ \alpha v_1+(-\alpha+\beta)v_2+(-\beta+\gamma)v_3+(-\gamma+\delta)v_4=0. \]
Since $v_1,v_2,v_3,v_4$ is linearly independent, we must have
\begin{align*}
0 &= \alpha,\\
0 &= -\alpha+\beta,\\
0 &= -\beta+\gamma,\\
0 &= -\gamma+\delta.
\end{align*}
From the first equation, we have $\alpha=0$. Substituting this into the second
equation gives us $\beta=0$. Substituting this into the third equation gives
us $\gamma=0$. Finally, substituting this into the fourth equation gives us
$\delta=0$. Thus, the only solution to the original equation is
$\alpha=\beta=\gamma=\delta=0$, and so the list is linearly independent.
\end{solution}

\begin{exercise}
Prove or give a counterexample: If $v_1,v_2,\dots,v_m$ is a linearly
independent list of vectors in $V$, then
\[ 5v_1-4v_2,\,v_2,\,v_3,\dots,v_m \]
is linearly independent.
\end{exercise}
\begin{solution}
Let $\alpha_1,\dots,\alpha_m\in\F$ such that
\[ \alpha_1(5v_1-4v_2)+\alpha_2v_2+\cdots+\alpha_mv_m=0. \]
Then we have
\[ 5\alpha_1v_1+(-4\alpha_1+\alpha_2)v_2+\alpha_3v_3+\cdots+\alpha_mv_m=0. \]
Since $v_1,v_2,\dots,v_m$ is linearly independent, we must have
\begin{align*}
0 &= 5\alpha_1,\\
0 &= -4\alpha_1+\alpha_2,\\
0 &= \alpha_3,\\
  &\vdotswithin{=}\\
0 &= \alpha_m.
\end{align*}
From the first equation, we have $\alpha_1=0$. Substituting this into the
second equation gives us $\alpha_2=0$. The remaining equations give us
$\alpha_k=0$ for $k=3,\dots,m$. Thus, the only solution to the original
equation is $\alpha_k=0$ for $k=1,\dots,m$, and so the list is linearly
independent.
\end{solution}

\begin{exercise}
Prove or give a counterexample: If $v_1,v_2,\dots,v_m$ is a linearly
independent list of vectors in $V$ and $\lambda\in\F$ with $\lambda\ne0$, then
$\lambda v_1,\lambda v_2,\dots,\lambda v_m$ is linearly independent.
\end{exercise}
\begin{solution}
Let $\alpha_1,\dots,\alpha_m\in\F$ such that
\[ \alpha_1(\lambda v_1)+\cdots+\alpha_m(\lambda v_m)=0. \]
Then we have
\[ \lambda(\alpha_1v_1+\cdots+\alpha_mv_m)=0. \]
Since $\lambda\ne0$, we must have
\[ \alpha_1v_1+\cdots+\alpha_mv_m=0. \]
Since $v_1,v_2,\dots,v_m$ is linearly independent, we must have $\alpha_k=0$
for $k=1,\dots,m$. Thus, the only solution to the original equation is
$\alpha_k=0$ for $k=1,\dots,m$, and so the list is linearly independent.
\end{solution}

\begin{exercise}
Prove or give a counterexample: If $v_1,\dots,v_m$ and $w_1,\dots,w_m$ are
linearly independent lists of vectors in $V$, then the list
$v_1+w_1,\dots,v_m+w_m$ is linearly independent.
\end{exercise}
\begin{solution}
Consider the following counterexample in $\R^2$:
\[ v_1=(1,0), \quad v_2=(0,1), \quad w_1=(0,1), \quad w_2=(1,0). \]
Then $v_1,v_2$ and $w_1,w_2$ are linearly independent lists of vectors in
$\R^2$. However,
\[ v_1+w_1=(1,1), \quad v_2+w_2=(1,1), \]
and so the list $v_1+w_1,v_2+w_2$ is linearly dependent. Thus, the statement
is false.
\end{solution}

\begin{exercise}
Suppose $v_1,\dots,v_m$ is linearly independent in $V$ and $w\in V$. Prove that
if $v_1+w,\dots,v_m+w$ is linearly dependent, then $w\in\spn(v_1,\dots,v_m)$.
\end{exercise}
\begin{solution}
Let $\alpha_1,\dots,\alpha_m\in\F$, not all zero, such that
\[ \alpha_1(v_1+w)+\cdots+\alpha_m(v_m+w)=0. \]
Then we have
\[ (\alpha_1v_1+\cdots+\alpha_mv_m)+(\alpha_1+\cdots+\alpha_m)w=0. \]
If $\alpha_1+\cdots+\alpha_m=0$, then we have
\[ \alpha_1v_1+\cdots+\alpha_mv_m=0, \]
which contradicts the linear independence of $v_1,\dots,v_m$. Thus, we must
have $\alpha_1+\cdots+\alpha_m\ne0$. Let $\beta=\alpha_1+\cdots+\alpha_m$.
Then we may write
\[ w = -\frac{\alpha_1}{\beta}v_1-\cdots-\frac{\alpha_m}{\beta}v_m. \]
Thus, $w$ is a linear combination of $v_1,\dots,v_m$, and so
$w\in\spn(v_1,\dots,v_m)$.
\end{solution}

\begin{exercise}
Suppose $v_1,\dots,v_m$ is linearly independent in $V$ and $w\in V$. Show that
\[ v_1,\dots,v_m,w \text{ is linearly independent} \iff
   w\notin\spn(v_1,\dots,v_m). \]
\end{exercise}
\begin{solution}
Suppose $v_1,\dots,v_m,w$ is linearly independent. If there were scalars
$\alpha_1,\dots,\alpha_m,\beta$ such that
\[ \alpha_1v_1+\cdots+\alpha_mv_m+\beta w=0, \]
then linear independence would imply that $\alpha_1=\cdots=\alpha_m=\beta=0$.
Suppose $w\in\spn(v_1,\dots,v_m)$. Then there exist scalars
$\gamma_1,\dots,\gamma_m$ such that
\[ \gamma_1v_1+\cdots+\gamma_mv_m=w. \]
Then we may write
\[ -\gamma_1v_1-\cdots-\gamma_mv_m+1\cdot w=0, \]
which contradicts the linear independence of $v_1,\dots,v_m,w$. Thus,
$w\notin\spn(v_1,\allowbreak\dots,\allowbreak v_m)$.

Conversely, suppose $w\notin\spn(v_1,\dots,v_m)$. Then there do not exist
scalars $\alpha_1,\dots,\alpha_m$ such that
\[ \alpha_1v_1+\cdots+\alpha_mv_m=w. \]
If $v_1,\dots,v_m,w$ is linearly dependent, then there exist scalars
$\beta_1,\dots,\beta_m,\gamma$, not all zero, such that
\[ \beta_1v_1+\cdots+\beta_mv_m+\gamma w=0. \]
If $\gamma=0$, then we have
\[ \beta_1v_1+\cdots+\beta_mv_m=0, \]
which contradicts the linear independence of $v_1,\dots,v_m$. Thus, we must
have $\gamma\ne0$. Then we may write
\[ w = -\frac{\beta_1}{\gamma}v_1-\cdots-\frac{\beta_m}{\gamma}v_m, \]
which contradicts the assumption that $w$ is not in the span of
$v_1,\dots,v_m$. Thus, $v_1,\dots,v_m,w$ is linearly independent.
\end{solution}

\begin{exercise}
Suppose $v_1,\dots,v_m$ is a list of vectors in $V$. For $k\in\{1,\dots,m\}$,
let
\[ w_k = v_1+\cdots+v_k. \]
Show that the list $v_1,\dots,v_m$ is linearly independent if and only if the
list $w_1,\dots,w_m$ is linearly independent.
\end{exercise}
\begin{solution}
Suppose $v_1,\dots,v_m$ is linearly independent. Let
$\alpha_1,\dots,\alpha_m\in\F$ such that
\[ \alpha_1w_1+\cdots+\alpha_mw_m=0. \]
Substituting $w_k=v_1+\cdots+v_k$ for each $k$, we get
\[ \alpha_1v_1+\alpha_2(v_1+v_2)+\alpha_3(v_1+v_2+v_3)+\cdots
   +\alpha_m(v_1+\cdots+v_m)=0. \]
Then we may rewrite this as
\begin{multline*}
(\alpha_1+\cdots+\alpha_m)v_1+(\alpha_2+\cdots+\alpha_m)v_2+\cdots\\
+(\alpha_{m-1}+\alpha_m)v_{m-1}+\alpha_mv_m=0.
\end{multline*}
Since $v_1,\dots,v_m$ is linearly independent, we must have
$\alpha_1=\cdots=\alpha_m=0$. Hence, $w_1,\dots,w_m$ is linearly independent.

Conversely, suppose $w_1,\dots,w_m$ is linearly independent. Let
$\beta_1,\dots,\beta_m\in\F$ such that
\[ \beta_1v_1+\cdots+\beta_mv_m=0. \]
Noting that $v_1=w_1$ and $v_k=w_k-w_{k-1}$ for $k\in\{2,\dots,m\}$, we may
rewrite the equation as
\begin{gather*}
\beta_1w_1+\beta_2(w_2-w_1)+\cdots+\beta_m(w_m-w_{m-1})=0,\\
(\beta_1-\beta_2)w_1+(\beta_2-\beta_3)w_2+\cdots+(\beta_{m-1}-\beta_m)w_{m-1}
+\beta_mw_m=0.
\end{gather*}
Since $w_1,\dots,w_m$ is linearly independent, we find that the system of
equations
\begin{align*}
0 &= \beta_1-\beta_2,\\
0 &= \beta_2-\beta_3,\\
  &\vdotswithin{=}\\
0 &= \beta_{m-1}-\beta_m\\
0 &= \beta_m,
\end{align*}
has the unique solution $\beta_1=\cdots=\beta_m=0$. Hence, $v_1,\dots,v_m$ is
linearly independent.
\end{solution}

\begin{exercise}
Explain why there does not exist a list of six polynomials that is linearly
independent in $\Poly_4(\F)$.
\end{exercise}
\begin{solution}
Since $\Poly_4(\F)$ can be spanned by the list $1,x,x^2,x^3,x^4$, and the
length of a linearly independent list cannot exceed the length of a spanning
list, there does not exist a list of six polynomials that is linearly
independent in $\Poly_4(\F)$.
\end{solution}

\begin{exercise}
Explain why no list of four polynomials spans $\Poly_4(\F)$.
\end{exercise}
\begin{solution}
Since the list $1,x,x^2,x^3,x^4$ is a linearly independent list in
$\Poly_4(\F)$ and has length $5$, and the length of a spanning list cannot be
less than the length of a linearly independent list, no list of four
polynomials can span $\Poly_4(\F)$.
\end{solution}

\begin{exercise}
Prove that $V$ is infinite-dimensional if and only if there is a sequence
$v_1,v_2,\dots$ of vectors in $V$ such that $v_1,\dots,v_m$ is linearly
independent for every positive integer $m$.
\end{exercise}
\begin{solution}
Suppose $V$ is infinite-dimensional. We proceed by induction to construct a
sequence of vectors $v_1,v_2,\dots$ in $V$ such that $v_1,\dots,v_m$ is
linearly independent for every positive integer $m$. Since $V$ is
infinite-dimensional, there exists a nonzero vector $v_1\in V$. Suppose we have
constructed a linearly independent list of vectors $v_1,\dots,v_m$ in $V$. If
every $v\in V$ is a linear combination of $v_1,\dots,v_m$, then
$V=\spn(v_1,\dots,v_m)$, which contradicts the assumption that $V$ is
infinite-dimensional. Thus, there exists a vector $v_{m+1}\in V$ that is not in
$\spn(v_1,\dots,v_m)$. Then by Exercise 2A13, the list
$v_1,\dots,v_m,v_{m+1}$ is linearly independent. By induction, we can construct
a sequence of vectors $v_1,v_2,\dots$ in $V$ such that $v_1,\dots,v_m$ is
linearly independent for every positive integer $m$.

Conversely, suppose there exists a sequence of vectors $v_1,v_2,\dots$ in $V$
such that $v_1,\dots,v_m$ is linearly independent for every positive integer
$m$. Then for every positive integer $m$, there exists a linearly independent
list of vectors in $V$ of length $m$. Since there is no maximum length for a
linearly independent list of vectors in $V$, this implies that $V$ is
infinite-dimensional.
\end{solution}

\begin{exercise}
Prove that $\F^\infty$ is infinite-dimensional.
\end{exercise}
\begin{solution}
By Exercise 2A17, it suffices to construct a sequence of vectors
$v_1,v_2,\dots$ in $\F^\infty$ such that $v_1,\dots,v_m$ is linearly
independent for every positive integer $m$. Let $v_k$ be the vector in
$\F^\infty$ whose $k$-th coordinate is $1$ and all other coordinates are $0$.
Explicitly, we have
\[ v_1=(1,0,0,\dots), \quad v_2=(0,1,0,\dots), \quad v_3=(0,0,1,\dots),
   \quad \dots \]
Then for any positive integer $m$, the list $v_1,\dots,v_m$ is linearly
independent, since the only way to write the zero vector as a linear
combination of $v_1,\dots,v_m$ is to take all coefficients to be zero. Thus,
by Exercise 2A17, $\F^\infty$ is infinite-dimensional.
\end{solution}

\begin{exercise}
Prove that the real vector space of all continuous real-valued functions on the
interval $[0,1]$ is infinite-dimensional.
\end{exercise}
\begin{solution}
By Exercise 2A17, it suffices to construct a sequence of continuous real-valued
functions $f_1,f_2,\dots$ on the interval $[0,1]$ such that $f_1,\dots,f_m$ is
linearly independent for every positive integer $m$. Let $f_k(x)=x^{k-1}$ for
each positive integer $k$. Then for any positive integer $m$, let
$\alpha_1,\dots,\alpha_m\in\R$ such that
\[ \alpha_1f_1(x)+\cdots+\alpha_mf_m(x)=0 \]
for all $x\in[0,1]$. This is equivalent to the polynomial equation
\[ \alpha_1+\alpha_2x+\cdots+\alpha_mx^{m-1}=0 \]
for all $x\in[0,1]$. Since a polynomial of degree at most $m-1$ can have at
most $m-1$ roots unless it is the zero polynomial, we must have
$\alpha_1=\cdots=\alpha_m=0$. Thus, the list $f_1,\dots,f_m$ is linearly
independent. Therefore, by Exercise 2A17, the real vector space of all
continuous real-valued functions on the interval $[0,1]$ is
infinite-dimensional.
\end{solution}

\begin{exercise}
Suppose $p_0,p_1,\dots,p_m$ are polynomials in $\Poly_m(\F)$ such that
$p_k(2)=0$ for each $k\in\{0,\dots,m\}$. Prove that $p_0,p_1,\dots,p_m$ is not
linearly independent in $\Poly_m(\F)$.
\end{exercise}
\begin{solution}
Since $p_k(2)=0$ for each $k\in\{0,\dots,m\}$, we have
\[ p_k(x)=(x-2)q_k(x) \]
for some polynomial $q_k(x)\in\Poly_{m-1}(\F)$. Then we may write
\[ p_0(x),p_1(x),\dots,p_m(x)=(x-2)q_0(x),(x-2)q_1(x),\dots,(x-2)q_m(x). \]
Since $\Poly_{m-1}(\F)$ can be spanned by the $m$ polynomials
$1,x,\dots,x^{m-1}$, the list of $m+1$ polynomials $q_0,q_1,\dots,q_m$ must be
linearly dependent. Thus, there exist scalars $\alpha_0,\dots,\alpha_m$, not
all zero, such that
\[ \alpha_0q_0(x)+\cdots+\alpha_mq_m(x)=0. \]
Then we may write
\[ \alpha_0p_0(x)+\cdots+\alpha_mp_m(x)
   =(x-2)(\alpha_0q_0(x)+\cdots+\alpha_mq_m(x))=0. \]
Since not all the scalars $\alpha_0,\dots,\alpha_m$ are zero, this shows that
$p_0,p_1,\dots,p_m$ is linearly dependent in $\Poly_m(\F)$.
\end{solution}
\end{exc}
